% 本文件由 LaTeX 排版 Agent 根据有效 translation.json 自动生成。
% author=Councils; work_id=3810; title=厄弗所大公会议（主后431年）

\chapter{厄弗所大公会议（主后431年）}

% structure: p0001 source=h1 target=omitted reason=page-title-already-rendered-by-parent

% structure: p0002 source=h2 target=section reason=relative-heading-rank; base=h2

\section{第一次会议}

% structure: p0003 source=h3 target=subsection reason=relative-heading-rank; base=h2

\subsection{会议记录选录}

\EditorialNote{在教宗代表抵达之前}\par

尼西亚会议陈述了这一信仰：\par

\Quote{我们信唯一的天主，等等。}\par

\Place{亚历山大里亚}的\Person{伯多禄}司铎，即\OriginalTerm{首席公证员}{primicerius}，说道：\par

我们手中持有最可敬的\Person{济利禄}总主教写给最可敬的\Person{聂斯多略}的书信，信中充满劝告与建言，因为他偏离了正确的信仰。如果您的圣洁（即神圣会议）下令，我将宣读此信……书信开头如下：\par

\OriginalTerm{有些人信口开河，据我所知，等等}{\GreekText{Καταφλυαροῦσι} \GreekText{μὲν}, \GreekText{ὡς} \GreekText{ἀκούω}, \GreekText{κ}.\GreekText{τ}.\GreekText{λ}}．\par

% structure: p0010 source=h3 target=subsection reason=relative-heading-rank; base=h2

\subsection{\Person{济利禄}致\Person{聂斯多略}书（\OriginalTerm{我听说有些人}{Intelligo quosdam meæ}）}

致最热心爱主、同为主仆的\Person{聂斯多略}，\Person{济利禄}在主内问安。\par

我听说有些人轻率地谈论我对您圣洁的看法，尤其常在当权者聚会的时候。也许他们以为这样说能讨您欢喜，但他们的话毫无意义，因为他们并没有受到我的不公正对待，反而被我揭露祇是为他们有益：这人欺压盲人和穷人，那人用剑伤害自己的母亲。另一个因为与侍女合谋偷窃他人的钱财，且一直背负着诸如此类无人愿加于最恶敌人身上的罪名。我不太计较这些人的言语，因为徒弟不能胜过师傅\BibleRef{玛 10:24}，我也不愿将我狭隘的脑筋伸展到教父之上，因为无论人走什么生活道路，都难以逃脱恶人的诽谤，他们的口满是咒骂和苦毒，最后必向众审判者交账。\par

但我要回到我特别思想的那一点。现在我劝你，作为在主内的兄弟，要极其精确地向民众提出教导的言语和信仰的教义，并且要考虑到：使一个信基督中的小子跌倒\BibleRef{玛 18:6}，会使身体遭受天主无法容忍的愤怒。当受害的群众如此众多时，需要何等勤勉和技巧，好让我们能向寻求的人施行真理的治愈之言。但我们将最出色地完成此事，如果我们不断翻阅圣教父的言语，并热切遵守他们的命令，考验自己是否按经上所记载的持守信仰，并使我们的思想符合他们正直无瑕的教导。\par

因此，神圣伟大的主教会议宣告：独生子，即按本性由天主圣父所生者，出自真天主的真天主，出自光的光，圣父藉祂创造了万有，祂降来，成了血肉，成了人，受难，第三日复活，升了天。我们当遵循这些话和这些法令，思量\Quote{天主圣言成了血肉、成了人}的意义。因为我们不说圣言的本性改变成了血肉，或转化成了由灵魂和肉身组成的完整的人；而是说圣言以位格的方式将具有理性灵魂的血肉与自己结合，以一种不可言喻、不可思议的方式成了人，并被称为人子；这并非仅仅因为祂愿意或被这样称呼，也不是因为祂取了一个位格，而是因为两性在真正的结合中相遇，由此产生了一个基督、一个圣子；因为两性的差异并未因结合而消失，反而神性和人性借着不可言喻、不可表达的合一，为我们造就了唯一的主耶稣基督。因此，那在万世之前就已存在、由圣父所生者，按血肉被说成是由一个女人所生，并非因为祂的神性在童贞女那里开始了存在——因为祂在圣父的出生之后不需要第二次出生（说那在万世之前就已存在、与圣父同永恒者需要第二次存在的开端，是荒谬和愚蠢的）——而是因为祂为了我们和我们的救恩，将人的身体与自己位格性地结合，从女人中出来，祂就这样按血肉被说成是出生的；因为祂并非先作为常人由童贞女所生，然后圣言降下进入祂里面，而是结合在母腹中就已经完成，祂被说成承受了按血肉的出生，将祂自己血肉的出生归给自己。为此，我们说祂受难并复活了；并非因为天主圣言在自己的本性中遭受了鞭打、钉子刺穿或任何其他创伤——因为神性是不能受苦的，因为它无形体——而是因为那成为祂自己身体的东西以这种方式受苦，祂也被说成为我们受难；因为祂本身不能受苦，却在一个受苦的身体中。同样，我们也以同样的方式理解祂的死亡；因为天主圣言按本性是不死不灭的，是生命和赋予生命的；然而，祂自己的身体，如保禄所说，因天主的恩宠为每一个人尝到了死味，祂自己就被说成为我们受死，并非因为祂在自己的本性中经历了死亡（这样说或想是疯狂的），而是因为，正如我刚才所说，祂的血肉尝到了死味。同样，祂的血肉复活了，这被说成是祂的复活，并非因为祂落入了朽坏（天主不许），而是因为祂自己的身体复活了。因此，我们宣认一个基督和主，并非将一个人与圣言一同崇拜（免得\Quote{与圣言一同}这个表达向思想暗示分裂），而是将祂作为同一个来崇拜，因为圣言的身体——祂以此与圣父同坐——并不与圣言本身分离，并非如同两个儿子与祂同坐，而是因与血肉的结合成为一位。然而，如果我们拒绝位格的结合，认为不可能或不合适，我们就落入谈论两个儿子的错误，因为那时将有必要加以区分，并说那真正是人的被尊为\Quote{子}的称号，而那真正是天主圣言的，按本性拥有\Quote{子}的名称和实质。因此，我们不可将唯一的主耶稣基督分为两个儿子。持有一种如同某些人所持的\Quote{位格的联合}也无助于正确的信德；因为圣经并没有说圣言将人的位格与自己联合，而是说祂成了血肉。然而，\Quote{圣言成了血肉}这个表达只能意味着：祂分享了与我们相同的血肉；祂使我们的身体成为自己的，从女人中出来成为人，并未抛弃祂作为天主的存在或祂由天主圣父所生，而是在取血肉时仍然保持祂所是。这正是正确信德的宣言所处处宣告的。这是圣教父们的信念；因此他们敢于称童贞女为天主之母，并非因为圣言的本性或神性起源于童贞女，而是因为从她而生的是那个有理性灵魂的神圣身体，圣言位格性地与其结合，被说成是按血肉而生的。因此，我如今因着基督的爱将这些事写给你们，如同弟兄恳求你们，并在基督和蒙选的天使面前劝勉你们：愿你们与我们一同思想和教导这些事，好使教会的和平得以保持，天主的司铎们之间的和好与爱德的纽带不断裂。\par

% structure: p0015 source=h3 target=subsection reason=relative-heading-rank; base=h2

\subsection{会议记录摘录（续）}

信读完后，亚历山大的主教济利禄说：这个神圣伟大的主教会议已听到我为捍卫正信而写给最虔诚的聂斯多略的信。我认为自己在任何方面都没有偏离信仰的正确陈述，即从前在尼西亚召集的神圣伟大主教会议所定的信经。因此，我恳请你们的圣善（即主教会议）说明：我所写的这些事是否正确、无过，并符合那次神圣主教会议。\par

\Emph{许多主教随后发表了意见，都赞同济利禄；在这些个人意见之后，会议记录继续（第491栏）。}\par

其余所有主教按各自等级依次作了同样的见证，并如此相信，正如教父们所陈述的，也正如最神圣的济利禄总主教给聂斯多略主教的信中所宣告的。\par

阿马西亚的主教帕拉迪乌斯说：接下来应当宣读最可敬的聂斯多略主教的信，就是那位最虔诚的司铎伯多禄曾提及的那封信；以便我们了解它是否与尼西亚教父的阐述一致。\par

在宣读了这封信后，亚历山大的主教济利禄说：对于刚刚诵读的信函，这神圣而伟大的大公会议有何看法？它是否也似乎与在尼西亚城集会的神圣会议所阐明的信仰相符？\par

\EditorialNote{主教们然后像先前一样，各自表达了意见，最后会议记录继续（col. 502）：}\par

所有主教一同高呼：凡不谴责聂斯多略者，当受绝罚。正信谴责这样的人；神圣的会议谴责这样的人。凡与聂斯多略共融者，当受绝罚！我们谴责聂斯多略的所有同党：我们一致谴责聂斯多略为异端者；凡与聂斯多略共融者，当受绝罚，等等，等等。\par

耶路撒冷的主教朱文纳尔说：让最神圣可敬的罗马教会总主教塞莱斯定关于信仰所写的信函被诵读。\par

\EditorialNote{塞莱斯定的信函被诵读，没有发表意见。}\par

亚历山大的司铎、公证首席官伯多禄说：与我们刚读到的一切完全一致的是我们最虔诚的主教济利禄所写的那些内容，我现在手头就有，若诸位虔诚者下令，我将诵读。\par

\EditorialNote{那封以此开头的信函被诵读：}\par

\OriginalTerm{\GreekText{Τοῦ} \GreekText{Σωτῆρος} \GreekText{ἡμῶν} \GreekText{λέγοντος} \GreekText{ἐναργῶς}, \GreekText{κ}.\GreekText{τ}.\GreekText{λ}.}{\GreekText{Τοῦ} \GreekText{Σωτῆρος} \GreekText{ἡμῶν} \GreekText{λέγοντος} \GreekText{ἐναργῶς}, \GreekText{κ}.\GreekText{τ}.\GreekText{λ}.}\par

% structure: p0028 source=h3 target=subsection reason=relative-heading-rank; base=h2

\subsection{济利禄致聂斯多略的\WorkTitle{救主明言}（\OriginalTerm{}{Cum salvator noster}）}

最可敬与虔敬的同事聂斯多略，济利禄及在亚历山大召开的埃及省会议，在主内问安。\par

当我们的救主明言说：\BibleQuote{爱父亲或母亲胜过我的人，不配属于我；爱儿子或女儿胜过我的人，不配属于我。}那么，我们这些您要求爱您胜过爱我们众人救主基督的人，将会怎样？在审判之日谁将帮助我们，或者我们能为如此长久地对您所行的亵渎保持沉默找到何种借口？倘若您独自受害，只是教导并持守这些事情，或许事小；但您已使整个教会大受绊倒，并将一种陌生而新奇的异端酵母投入民众之中。不止是那里（即君士坦丁堡），而且遍及各处的（您的解释著作）都被寄送。在这种情况下，我们如何能再为我们的沉默辩护，或者我们怎能不被迫记起基督曾说：\BibleQuote{不要以为我来是给地上送和平：我来不是送和平，而是刀剑。因为我来是为使人与父亲对立，女儿与母亲对立。}因为如果信仰受损，就让对父母的尊重像陈腐摇晃之物般失去；甚至对儿女和兄弟的柔情之律也要静默；对虔诚者来说，死比活着更好，\BibleQuote{为要得着更美的复活}，如经上所载。\par

看，因此，我们与在大会罗马集会的神圣会议，由其主席最神圣可敬的兄弟和同事、主教塞莱斯定，也通过这第三封信向你作证，并劝诫你放弃你所持守和教导的这些有害而歪曲的教义，而接受从起初通过圣宗徒和传福音者传给教会的正信，他们\BibleQuote{从起初见过和做过圣言仆役}。如果你的虔敬不按照我们有福记忆的兄弟、最可敬的同事、罗马教会主教塞莱斯定信函中所定的界限而达此意，那么请确信你在天主面前与我们无分，也无地位或名分（\OriginalTerm{\GreekText{λόγον}}{\GreekText{λόγον}}）在天主的司铎和主教中。因为我们不可能忽视如此被扰乱的教会、被绊倒的民众、被废弃的正信以及被你——本应拯救他们的人——所驱散的羊群，如果我们自己确实是正信的持守者、圣父们虔诚的追随者。我们与所有那些因信仰被你的虔敬逐出或罢黜的平信徒和神职人员共融；因为那些决心正确信从的人不应因你的选择而受苦；因为他们反对你做得对。这一点你已在致我们至圣的共主教罗马的塞莱斯定的信函中提及。\par

但你的敬意仅仅与我们一道承认从前在尼西亚召开的大公神圣会议由圣神所颁布的信仰象征书是不够的：因为你没有正确地持守和解释它，而是相反地歪曲了；尽管你用言语承认文字形式。但此外，你必须在书写和起誓中，承认你也谴责你自己那些污秽不敬的教义，并且你将持守并教导我们——东西方的所有主教、教师和民众领袖——所持守的。罗马的神圣会议和我们全体都同意从亚历山大教会写给您的虔敬的信函，认为是正确无误的。我们附加了我们自己的信函，以及你必须持守和教导的内容，以及你应当小心避免的内容。这就是大公教会和宗徒教会的信仰，所有正统的主教，无论东方西方，都同意：\par

\Quote{我们信唯一的天主，全能的圣父，造天地万物有形无形者；并信唯一的主耶稣基督，天主的独生子，在万世之前由圣父所生，即出自圣父的本体：天主出自天主，光明出自光明，真天主出自真天主，受生而非受造，与圣父同性同体；万有借着他而造，无论天上地上万物。他为了我们人类并为了我们的救恩，从天降下，取得肉体，成为人。他受难，第三日复活。他升天，从那里将再来审判生者死者。并信圣神。但那些说：曾有一段时间他不在，并被生之前他不存在，或者说他出于先前不存在之物，或他出于另一种实体或本质，以及天主子是可变化或改变的——这样的人，大公教会和宗徒教会予以绝罚。}\par

我们完全追随圣教父们所作的宣认（在他们内发言的圣神所作的），追随他们意见的宗旨，如同走在王道之上，我们宣认：天主的独生圣言，由圣父同一本体所生，真天主出自真天主，光出自光，藉着祂万有得以造成，天上的和地上的，祂为我们的救恩降下，\OriginalTerm{空虚自己}{\GreekText{καθεὶς} \GreekText{ἑαυτὸν} \GreekText{εἰς} \GreekText{κένωσιν}}，降生成人并成为人；即从\Person{童贞玛利亚}取得肉体，并从母腹中使之成为自己的，祂为我们而受生，从女人生出为人，并未抛弃祂所是的；但祂虽取了血肉，仍保持祂所是的，在本质和真理上是天主。我们也不说祂的肉体变成了神性的本性，也不说天主圣言那不可言喻的本性被搁置以取血肉的本性；因为祂是永不改变的、绝对不变易的，常是同一的，按照圣经。因为祂虽可见、是婴孩，包裹在襁褓中，甚至在\Person{童贞玛利亚}的怀中，仍作为天主充满万有，并与生祂的父同掌权，因为天主性是无量度、无维度的，不能有限制。\par

宣认圣言按照本体与肉体成为一体，我们钦崇一位圣子和主耶稣基督：我们不将天主与人分开，也不将祂分成部分，好像两个本性在祂内仅仅通过尊荣和权威的分享而彼此结合（因为那是新奇之事，绝非其它），我们也不分别将\Quote{基督}这名字给予天主圣言，又将同一名字分别给予一个不同的从女人所生的；但只知道一位基督，即来自圣父天主的圣言，带着祂自己的肉体。因为作为人，祂与我们一同受傅，虽然正是祂自己将圣神赐予那些堪当的人，且没有限量，按照真福\Person{圣史若望}的话。\par

但我们不说天主圣言住在他内如同在一个由\Person{童贞玛利亚}所生的普通人内，免得基督被想成一位承载天主的人。因为虽然圣言在我们中间支搭帐篷，但经上又说：\BibleQuote{有形地居住一切天主性之圆满}；但我们理解祂成了血肉，并不只是如同经上说居住在圣人们内那样，但我们界定，祂内支搭帐篷的方式是按照同等（\OriginalTerm{按照同样的方式}{\GreekText{κατὰ} \GreekText{τον} \GreekText{ἴσον} \GreekText{ἐν} \GreekText{αὐτῷ} \GreekText{τρόπον}}）。但祂被造成合一按照本性（\OriginalTerm{按照本性}{\GreekText{κατὰ} \GreekText{φύσιν}}），而非变成血肉，祂的居住方式是这样的，就如同我们说人的灵魂在自己身体内那样。\par

所以基督是圣子也是主，并非是一个人仅仅达到了与天主这样一种结合，即在于单一的尊荣或权威的合一。因为团结本性的不是荣誉的相等；因为若如是，则\Person{伯多禄}和\Person{若望}，他们彼此有相等的荣誉，同为宗徒和圣门徒，［可能成为一体，］但二者并非一体。我们也不理解结合的方式是并置，因为这不满足于\OriginalTerm{自然的合一}{\GreekText{πρὸς} \GreekText{ἕνωσον} \GreekText{φυσικήν}}。也不是按照相对的参与，就如我们也与主结合，正如经上记载：\BibleQuote{我们与祂成为一灵}。反之，我们不赞成\Quote{\OriginalTerm{结合}{\GreekText{συναφείας}}}一词，因为它没有充分表示合一。但我们不称来自圣父的天主圣言为\Quote{基督的天主}或\Quote{基督的主}，免得我们公开地将一位基督、圣子和主分割为二，并陷于亵渎的指控，使祂成为自己的天主和主。因为天主圣言，如我们已经说过的，在肉体中成为本位的合一，但祂是万有之神并掌管万有；但祂不是自己的奴仆，也不是自己的主。因为如此想或如此教是愚蠢的，甚至是不敬的。因为祂说天主是祂的父亲，虽然祂按本性是天主，并出自祂的本体。然而我们并非不知道，当祂保持为天主时，祂也成了人并服从天主，按照适合人性本性的法律。但祂如何能成为自己的天主或主呢？所以，作为人，并考虑到祂屈尊的尺度，祂被说成是与我们同等服从天主的；因此祂成了属于法律的，虽然作为天主祂颁布法律并是立法者。\par

我们也小心如何论及基督：\Quote{我因那穿衣者而崇拜那被穿衣者，并因那看不见者而崇拜那可见者。}与此相关的，这样说也是可怕的：\Quote{被取的和取用的都具有天主之名。}因为这种说法再次将基督分为二，并把人和天主分别各自独立。因为这种说法公开否认合一，按照这合一，一个不在另一个内被崇拜，天主也不与另一个并存；但耶稣基督被视为一位，独生子，应连同祂自己的肉体一起以同一朝拜受尊敬。\par

我们宣认祂是圣子，由圣父天主所生，独生的天主；虽然按照祂自己的本性，祂不受苦，但祂按照圣经为我们在肉体中受苦，虽不受苦，但在祂被钉的身体中，祂使自己承受了自己肉体的苦难；因天主的恩宠，祂为众人尝了死味：祂为此交付了自己的身体，虽然祂按本性就是生命和复活，为的是藉祂不可言喻的大能踏平死亡，首先在自己肉体中，成为死者中的首生者，睡了之人的初果。并且为使人的本性获得不朽，因天主的恩宠（如我们刚才所说），祂为每个人尝了死味，三天后复活，劫掠了地狱。所以，虽然经上说是藉着人死人复活，但我们理解那人是天主圣言，死亡的权势藉祂而被解除，祂将在时期圆满时作为一位圣子和主，在圣父的荣耀中来临，为按公义审判世界，正如经上所记载的。\par

因为当他作为天主谈论自己时，说：\BibleQuote{谁看见了我，就看见了父}，以及\BibleQuote{我与父原是一体}，我们便默观他那不可言喻的天主性本质，藉此他因本质的同一而与父为一——\BibleQuote{他是他光荣的形象、印记和光辉。}但当他不轻看自己人性的限度时，他对犹太人说：\BibleQuote{如今你们却想杀害我，这个把真理告诉你们的人。}同样，我们毫不逊色地从前者认出他是天主的圣言，因他与父的同一和相似，以及他人性的境况。因为如果必须相信，那按本性是天主的，成了血肉，即成为具有理性灵魂的人，那么某些人对这些合适于他人性的话感到羞愧，又有什么理由呢？因为如果他拒绝那些合适他人性的话，谁又迫使他成为像我们一样的人呢？正如他为我们自卑，甘愿自我空虚（\OriginalTerm{自我空虚}{\GreekText{κένωσιν}}），那么有什么理由有人会拒绝那些合适于此空虚的话呢？因此，凡福音中所读到的一切话，都应归于一个位格，即降生成人的圣言的一个位格。因为主耶稣基督是同一的，按照圣经，尽管他被称为\BibleQuote{我们所宣认的宗徒和大司祭}，向天主和父奉献我们对他所宣认的信德，并藉着他向天主、父及圣神奉献；但我们说他按本性是天主的独生子。我们并不将司祭职的名称和事实归于任何与他不同的人，因为他成了\BibleQuote{天主与人之间的中保}，和平的调解者，将自己作为馨香的祭献献给天主和父。因此他也说：\BibleQuote{牺牲和素祭，你不喜欢，却为我预备了身体；全燔祭和赎罪祭，你不喜爱。于是我说，看，我来了（书卷上已记载我），为承行你的旨意，天主。}因为为了我们，他献上了自己的身体作为馨香的祭献，而非为自己；因为为自己需要什么祭献或牺牲呢？他作为天主，超越一切罪过。因为\BibleQuote{所有人都犯了罪，失掉了天主的光荣}，以致我们易于跌倒，人的本性陷入了罪恶，但他却不然（因此我们失掉了他的光荣）。那么，真羔羊为我们、因我们的缘故而死，还能有怀疑吗？而说他是为自己和我们而奉献自己，则绝不能免于不虔敬的指控。因为他从未犯过任何过失，也没有犯罪。那么他需要什么祭献呢？他没有罪，而祭献正是为罪而设的。但当他论及圣神时，他说：\BibleQuote{他要光荣我。}若我们正确思考，我们并不说那同一的基督和圣子因需要从另一处得光荣而从圣神领受了光荣；因为圣神既不比他大，也不在他之上，而是因为他藉圣神在奇事中彰显自己的天主性，故他说自己被他光荣了，就好像我们中任何人论及自己内在的能力，例如，或对某事的认识，说\Quote{他们光荣了我。}因为尽管圣神是同一本质，但我们独自默想他，作为圣神而非圣子；但他并不与他不同；因为他被称为真理之神，而基督是真理，且由他所派遣，正如他也是由天主和父所派遣。因此，当我们的主耶稣基督升天之后，圣神藉圣宗徒的手行奇迹时，他便光荣了他。因为人们相信那藉自己的圣神行事者，按本性是天主。因此他说：\BibleQuote{他要从我领受，并告诉你们。}但我们这样说，并非因为圣神是藉着与另一位的分享而有智慧和能力；因为他全备，一无所缺。因此，既然他是父（即圣子）的能力和智慧的圣神，他显然是智慧和能力。\par

既然圣童贞女按肉体生下了那按本性成为与肉体合一的天主，因此我们也称她为天主之母，并非因为圣言的本性从肉体开始存在。\par

因为\BibleQuote{在起初已有圣言，圣言与天主同在，圣言就是天主}，他是万世的创造者，与父同永，万有的创造者；但如我们已说过的，既然他使自己位格性地与她子宫中的人性结合，他也使自己作为人而出生，并非因他的本性需要在时间中以及在世代的末期中出生，而是为祝福我们存在的开端，并为使那将我们全族的肉体推向死亡的力量，因他由女人在肉身中出生而失去未来的能力。并且这话：\BibleQuote{你在痛苦中要生育子女}，因他而被除去，他显示了先知所言的真确：\BibleQuote{强力的死亡吞没了他们，但天主又擦去了各人脸上的眼泪。}因此我们也说他被邀请参加了加里肋亚加纳的婚宴，并祝福了它，连同他的圣宗徒，依照救恩计划。我们由圣宗徒和圣史们，以及所有天主所启示的圣经，以及真福教父们的真实宣认，领受了这些。\par

对于这一切，你的尊威也应同意并注意，毫无欺诈。至于你的尊威必须绝罚之事，我们已附于本函之后。\par

% structure: p0044 source=h3 target=subsection reason=relative-heading-rank; base=h2

\subsection{圣济利禄驳聂斯多略的十二绝罚}

若有人不承认厄玛奴耳是真天主，因此圣童贞女是天主之母（\OriginalTerm{天主之母}{\GreekText{Θεοτόκος}}），因为她在肉体中生了降生成人的天主圣言【如经上记载：\BibleQuote{圣言成了血肉}】，此人当受绝罚。\par

\Quote{聂斯多略：若有人说厄玛奴耳是真天主，而非天主与我们同在，即祂将与我们相似的本性（祂从童贞玛利亚所取的）与自己结合，并居住其中；又若有人称玛利亚为天主圣言之母，而非称她为厄玛奴耳之母；并主张天主圣言改变成祂仅为使祂的神性可见而取的肉体，并以人的形像被找到——此人当受绝罚。}\par

若有人不承认天主父之圣言与肉体是本体结合，且同一基督同时是天主亦是人的肉体是祂自己的，而一并成为唯一的基督——此人当受绝罚。\par

\Quote{聂斯多略：若有人在圣言与肉体的结合上，主张神圣本质从一个地方移到另一个地方；或说肉体能够接纳神圣本性，且部分地与肉体结合；或由于肉体接纳天主而将之归于无限无界的延伸，并说天主和人本性相同——此人当受绝罚。}\par

若有人在本体结合之后，在唯一的基督内分割诸本体，仅凭尊威或权能与能力那样的关联加以结合，而非因（\OriginalTerm{结合}{\GreekText{συν}\'o\GreekText{δῳ}}）而成为（\OriginalTerm{本性结合}{\'\GreekText{ενωσιν} \GreekText{φυσικ}\'\GreekText{ην}}）——此人当受绝罚。\par

\Quote{聂斯多略：}若有人说基督（也就是厄玛奴耳）是唯一的，并非因关联而是因本性，且不承认两性（圣言与被取的人性）在一位圣子内是（\OriginalTerm{关联}{\GreekText{συν}\'\GreekText{αφεια}}）而仍不混杂——此人当受绝罚。\par

若有人将福音书和宗徒书信中、或圣人们所论及基督的、或祂自己所表达的（\OriginalTerm{用语}{\GreekText{φων}\'\GreekText{ας}}）分给两个位格或自立体，并将其一些用于那个与天主圣言分离的人，另一些用于唯独的天主父之圣言，且认为适用于天主——此人当受绝罚。\par

\Quote{聂斯多略：若有人将福音与宗徒书信中论及基督两性的用语只归于其中一性，并将受苦归于神圣圣言，既在肉体也在神性内——此人当受绝罚。}\par

若有人胆敢说基督是位具有天主者，而非是本性上的独生子真天主，因为\Quote{圣言成了血肉}且\Quote{祂也同样分享了血肉}——此人当受绝罚。\par

\Quote{聂斯多略：}若有人胆敢主张在取了人性后，仍只有一位本性上的圣子（\OriginalTerm{本性之子}{naturaliter filius}，即圣言），而祂在取了肉体后无疑是厄玛奴耳——此人当受绝罚。\par

若有人胆敢说天主父之圣言是基督的天主或基督的主，而不承认祂同时是天主和人，因为经上记载：\Quote{圣言成了血肉}——此人当受绝罚。\par

\Quote{聂斯多略：\Quote{若有人在降生成人后，称基督之外的另一个为圣言，并胆敢说仆人的形像与天主圣言同等无始无受造，而非由祂作为自己本性的主、创造主和天主所造，且祂曾应许再建它说：\InnerQuote{拆毁这座圣殿，三天之内我要把它重建起来}——此人当受绝罚。}}\par

若有人说耶稣作为人，只是被天主圣言所推动，而祂之享有独生子的光荣，并非祂所固有——此人当受绝罚。\par

\Quote{聂斯多略：}若有人说那由童贞女所形成的人是独生子，生于圣父怀前，在晨星之前，\BibleRef{咏 108:3}，而不承认祂因与那本性为父之独生子的关联而获得\Quote{独生子}的称号；此外，若有人称基督为厄玛奴耳之外的另一个——此人当受绝罚。\par

若有人胆敢说那被取的人（\OriginalTerm{被取之人}{\'\GreekText{αναληφθ}\'\GreekText{εντα}}）应与天主圣言同受崇拜、同受光荣、同被承认为天主，却是两件不同的事，一个与另一个（因为\Quote{同}字被添加以表达此意），而不以同一个崇拜敬拜厄玛奴耳并归予同一光荣，正如经上所写：\Quote{圣言成了血肉}——此人当受绝罚。\par

\Quote{聂斯多略：若有人说仆人的形像应因其本身，即就其本性的关系，受尊敬，且是万物的统治者，而非因其与那独生子的神圣、本身统御万有的本性之关联而受尊敬——此人当受绝罚。}\par

若有人说唯一的主耶稣基督靠圣神受光荣，以致祂使用了并非自己的权能，并从祂那里领受了制伏不洁鬼神的力量和在世人面前行奇迹的能力，而不承认祂借着自己的圣神行了这些神迹——此人当受绝罚。\par

\Quote{聂斯多略：若有人说仆人的形像与圣神同本性，而非承认它从成胎起就因圣言的中介而与圣言结合，并借此在人间施行奇迹的治愈，且有驱逐魔鬼的能力——此人当受绝罚。}\par

凡说那成了血肉、成为像我们一样的人，并不是天主圣言本身，而是另一个人，一个妇人所生的人，却与祂不同（\OriginalTerm{截然不同的人}{\GreekText{ἰδικῶς} \GreekText{ἄνθρωπον}}），且成了我们的大司祭和宗徒；或若有人说，祂是为自己而非为我们献上了自己作为祭献，而祂本无罪，无需奉献或祭献：此人当受绝罚。\par

\Quote{聂斯多略：若有人主张那从起初就有的圣言，成了我们所宣认的大司祭和宗徒，并为我们献上了自己，而不说作宗徒乃是厄玛奴耳的工程；又若有人如此分割那结合者（圣言）与被结合者（人性）之间的祭献，将其归于一个共同的子嗣，即不将属于天主的归于天主，属于人的归于人：此人当受绝罚。}\par

凡不承认主的肉是赋予生命的，且这肉属于天主圣父的圣言，为祂所固有，反而妄称这肉属于另一个与祂（即圣言）仅在尊荣上结合、且作为神性居所的人；且不按我们所说的承认那肉赋予生命，因为它是那赋予万有生命的圣言之肉：此人当受绝罚。\par

\Quote{聂斯多略：若有人主张那与天主圣言结合的肉，凭其本性的能力是赋予生命的，而主自己却说：\BibleQuote{使人生活的是神，肉一无所益}（\BibleRef{若 6:61}），此人当受绝罚。[他又说：\BibleQuote{天主是神}（\BibleRef{若 4:24}）。因此，若有人主张天主圣言在肉身方面、在其本体上成了血肉，并坚持这样理解主基督；而主在复活后对门徒说：\BibleQuote{你们摸摸我，看看；因为神没有肉和骨头，如同你们看见我有的}（\BibleRef{路 24:39}）：此人当受绝罚。]}\par

凡不承认天主圣言在肉身上受了苦，在肉身上被钉十字架，又同样在肉身上尝了死味，且成了死者中的首生者，因为祂是天主，是生命，是祂赋予生命：此人当受绝罚。\par

\Quote{聂斯多略：若有人在宣认肉身的苦难时，也把这些归于天主圣言，如同归于祂所显现的肉身，从而不区别各本性的尊荣：此人当受绝罚。}\par

% structure: p0069 source=h3 target=subsection reason=relative-heading-rank; base=h2

\subsection{会议记录摘录（续）}

\EditorialNote{[会议记录未记载采取了任何行动。有人口头报告说，过去三天见过聂斯多略的人，对任何悔改已不抱希望。经斐立比主教弗拉维安动议，宣读了几段教父著作；随后又读了聂斯多略著作中的一些选段。接着宣读了迦太基总主教卡普雷奥卢斯的信，为其缺席致歉；该信并未直接提及聂斯多略，而是恳请会议确保不容忍任何新奇事物。读信后，会议记录继续（第534栏）。]}\par

亚历山大里亚教会主教济利禄说：最可敬和虔诚的迦太基主教卡普雷奥卢斯这封已宣读的信，包含了最清晰的见解，应将其载入会议记录。因为它希望确认信仰的古老教义，并谴责和禁止那些荒谬设想、不虔敬产出的新奇事物。\par

所有主教同时喊道：这是我们众人（\OriginalTerm{声音}{\GreekText{φωναί}}）的呼声，这是我们大家所说的话——完成此事是我们众人的愿望。\par

\EditorialNote{[紧接着是罢免判决和签名。几乎可以肯定，此处有遗漏，很可能遗漏了关于济利禄《十二绝罚》的整个讨论。]}\par

% structure: p0074 source=h3 target=subsection reason=relative-heading-rank; base=h2

\subsection{大公会议反对聂斯多略的法令}

\Emph{（见于所有希腊文及拉丁文译本 \WorkTitle{公会议汇编} 中。）}\par

鉴于不虔敬的聂斯多略，除其他事项外，未服从我们的传唤，也未接待我们派去的圣主教们，我们被迫审查他的不敬神学说。我们发现他在书信、论文以及在本城所作的、已有证词的演讲中，持有并传播不敬神的学说。受教规和我们的最圣父、同为仆人的罗马主教策肋定之信（\OriginalTerm{迫于教规和信函等}{\GreekText{ἀναγκαίως} \GreekText{κατεπειχθέντες} \GreekText{ἀπό} \GreekText{τε} \GreekText{τῶν} \GreekText{κανόνων}, \GreekText{καὶ} \GreekText{ἐκ} \GreekText{τὴς} \GreekText{ἐπιστολῆς}, \GreekText{κ}.\GreekText{τ}.\GreekText{λ}}）所迫，我们含着许多眼泪，对他宣布这一悲伤的判决：即被祂亵渎的我们的主耶稣基督，藉此神圣会议决定，聂斯多略被排除于主教尊位及一切司祭共融之外。\par

% structure: p0077 source=h2 target=section reason=relative-heading-rank; base=h2

\section{第二次会议}

% structure: p0078 source=h3 target=subsection reason=relative-heading-rank; base=h2

\subsection{会议记录摘录}

最虔诚、最敬爱天的主教阿卡狄乌斯和普洛耶克图斯，以及最蒙天主钟爱的腓理伯司铎、宗座代表，随即进入就座。\par

司铎、宗座代表腓理伯说：我们赞颂圣而可钦的天主圣三，因我们的卑微堪当参加你们的圣会议。因为许久以前（\OriginalTerm{许久以前}{\GreekText{πάλαι}}），我们最圣洁、有福的教宗策肋定，宗座主教，通过致亚历山大里亚主教、圣洁且至虔诚者济利禄的信函，已对本案和事务作出了判决（\OriginalTerm{已判决}{\GreekText{ὥρισεν}}），这些信函已向你们圣会展示。现在，为巩固大公（\OriginalTerm{大公}{\GreekText{καθολικῆς}}）信仰，他又通过我们致信各位圣善者，请你们命令（\OriginalTerm{命令}{\GreekText{κελούσατε}}）恭敬有加地（\OriginalTerm{合宜地}{\GreekText{πρεπόντως}}）宣读，并录入教会记录中。\par

罗马教会主教代表、主教阿卡狄乌斯说：请诸位有福者下令，宣读我们带来的圣洁、永受敬仰的教宗策肋定、宗座主教的书信，从这些信中，你们的可敬者将能看出他对所有教会何等关怀。\par

普洛耶克图斯，一位主教和罗马教会的代表说：请予允准等等。[与阿尔卡狄乌斯所说\Emph{完全相同}，逐字!]\par

随后，至圣及钦崇天主的亚历山大教会主教济利禄按顺序发言如下；罗马圣\OriginalTerm{公教会}{\GreekText{καθολικῆς}}的公证人西利西乌斯宣读了它。\par

亚历山大主教济利禄说：请以合宜的敬意，向神圣主教会议宣读由至圣及全然蒙福的罗马宗座主教塞莱斯定所收的信函。\par

罗马城圣\OriginalTerm{公教会}{\GreekText{καθολικῆς}}的公证人西利西乌斯宣读了它。\par

宣读拉丁文后，耶路撒冷主教犹文纳尔说：请将刚刚宣读的至圣蒙福的伟大罗马主教的文书载入会议记录。\par

所有极其可敬的主教们请求将信函翻译并宣读。\par

宗座长老及代表斐理伯说：惯例已充分遵守，即宗座的文书应先以拉丁文宣读。但现在既然您的圣洁要求也以希腊文宣读，必须满足您的愿望；我们已确保此事完成，并将拉丁文译成希腊文。因此请下令，使它在您神圣的聆讯中被接受和宣读。\par

主教及代表阿尔卡狄乌斯和普洛耶克图斯说：正如您的福乐所命令，应将我们带来的文书告知众人，因为我们神圣的弟兄主教中有不少不懂拉丁文，因此信函已译成希腊文，若您下令，请予宣读。\par

斐立伯主教弗拉维安说：请接受并宣读至圣钦崇天主的罗马教会主教信函的译文。\par

亚历山大的长老及首席公证人伯多禄宣读如下：\par

% structure: p0092 source=h3 target=subsection reason=relative-heading-rank; base=h2

\subsection{教宗塞莱斯定致以弗所主教会议的书信}

塞莱斯定主教致在厄弗所集会的圣主教会议，可敬爱慕的弟兄们，愿在主内安好。\par

司祭会议见证圣神的临在。因我们所读的是真实的，真理不能撒谎，即福音的应许：\BibleQuote{哪里有两个或三个人，因我的名字聚在一起，我就在他们中间。}既然是这样，若圣神不缺席于如此小的数目，那么当如此众多的圣者聚集时，我们岂不更可相信祂临在吗？每一次主教会议因其应得的特殊尊敬而神圣；因为在每一次这样的会议中，都要顾及对那宗徒的最著名会议的敬仰，我们读到过这会议。他们的导师，即他们所领受传报的主，从未缺席于他们，而常常以主和导师的身份临在；那些教导的人也从没有被他们的老师遗弃。因为派遣他们的就是他们的老师；命令应教导什么的，就是他们的老师；祂宣称在祂的宗徒中自己被听到的，就是他们的老师。这传扬的职责已委托给所有主的司祭共同承担，因为凭继承之权我们有义务承担这份挂虑，无论我们中的谁在各处代替他们宣扬主的名，因为祂对他们说：\BibleQuote{你们去，教训万民。}亲爱的弟兄，你们应注意，我们领受了一个普遍的诫命：因为祂愿意我们所有人都履行祂如此共同托付给所有宗徒的职务。我们必要追随我们的先辈。那么，让我们所有人都承担他们的劳苦，因为我们继承他们的尊位。我们藉着宣讲他们所教导的同一教义来显明我们的勤奋，除此以外，按照宗徒的劝诫，我们被禁止添加任何东西。因为保管所交托给我们的宝藏的职分，并不比传递它的职分逊色。\par

他们撒下了信德的种子。这将是我们的挂虑：愿我们伟大家主（无疑这宗徒的圆满唯独归于祂）的来临，会找到不腐而多产的果实。因为蒙选器皿告诉我们，种植和浇灌是不够的，除非天主赐予增长。因此我们必须共同努力，保持那藉由宗徒继承而传到今日的信德。因为我们期望按宗徒行事。因为现在受质疑的不是我们的\Emph{外表}，而是我们的信德。我们必须拿起属神的武器，因为战争是思想之战，武器是话语；这样我们将在我们君王的信德中强壮。蒙福的宗徒保禄劝勉众人应留在那地方，他命令弟茂德留在那里。同一地方，同一缘由，交给我们同一职责。现在让我们也做并研习他当时命令他所做的。且勿让人有别的心思，勿让人留意过份离奇的故事，正如他亲自命令的。让我们同心合意，思念同一件事，因为这是有益的；让我们做事不凭竞争，不凭虚荣；让我们在一切事上同心同意，当唯一的信德受到攻击时。愿整个身体与我们一同哀伤同悲。那位将要审判世界的被传唤受审；那位将要批判所有的，自己成了批判的对象；那位救赎我们的，遭受诽谤。亲爱的弟兄，你们要以天主的武装束身。你们知道什么盔甲要保护我们的头，什么铠甲保护我们的胸。因为这不是教会营垒第一次接纳你们作其统帅。无人怀疑，借着使二者合为一的主的恩宠，将有和平，武器将被放下，因为起因本身为自己辩护。\par

让我们再看一次我们导师的这些话，祂特意针对主教说：\BibleQuote{你们要留心自己和整个羊群，圣神立你们为主教，管理天主的教会，就是祂用自己的血所获得的。}\par

我们读到那些在厄弗所（正是您的圣座聚集的地方）听到这事的人，是从那里被召叫的。因此，对于那些知道这信仰宣讲的人，也让您对这同一信仰的辩护为他们所知。让我们以对重大事务应有的敬畏向他们展示我们心灵的坚毅；这些事长久以来由和平以虔诚的理解所守护。\par

让您宣布那些从宗徒们完好保存下来的事；因为专制反对的话语绝不许可对着万王之王提出，真理的事业也不能被谎言压制。\par

我最蒙福的弟兄们，我劝勉你们：唯当珍视爱德，我们应在爱德中持守，正如宗徒若望的声音所说，我们在此城敬礼其圣髑。让我们向主献上共同的祈祷。因为藉由思考如此众多司铎联合祈祷时天主临在的能力，我们可以形成某种概念，正如我们读到，十二位曾一起恳求的地方本身也被感动。而宗徒们那次祈祷的用意是什么？是为使他们获得恩宠，得以放胆宣讲天主圣言，并藉其大能行事；两者都因我们的天主基督的惠赐而获得。现在，除了使你们能放胆宣讲上主的圣言之外，你们的圣洁会议还要求什么呢？还有什么比求祂赐予你们恩宠以保存祂所托付你们宣讲的更为重要呢？这样，你们充满圣神，正如经上所载，虽以不同的声音，却能阐述那同一的真理，即圣神亲自教导你们的。\par

简而言之，为这一切所激励（因为，正如宗徒所说：\Quote{我向那认识法律的人说话，且在成全的人中讲论智慧}），请你们坚守大公信仰，捍卫各教会的和平，因为如此说，既是向过去、现在，也是向将来，祈求并保守\Quote{那些属于耶路撒冷和平的事}。\par

出于我们的关怀，我们派遣了与我们同心且备受推崇的圣弟兄和同僚司铎——主教阿契迪乌斯、普洛耶克图斯，以及我们的司铎斐理伯——让他们临场所行之事，并执行那些已经由我们先前所规定的事（\OriginalTerm{已经由我们先前所规定的事}{quæ a nobis antea statuta sunt, exequantur}）。\par

对于执行此事，我们毫不怀疑您的圣座将会同意，因为所规定的事是为整个教会的安全。发布于五月八日，在巴苏斯和安提阿库斯担任执政官之年。\par

% structure: p0103 source=h3 target=subsection reason=relative-heading-rank; base=h2

\subsection{从会议记录摘录}

同时，所有最可敬的主教们呼喊说：这判决是正义的！致色勒斯丁，一位新保禄！致济利禄，一位新保禄！致色勒斯丁，信仰的守护者！致色勒斯丁，与会议同心者！致色勒斯丁，整个会议表示感谢！唯一的色勒斯丁！唯一的济利禄！会议唯一的信仰！世界唯一的信仰！\par

普洛耶克图斯，最可敬的主教及代表说：愿您的圣座考虑圣可敬教宗色勒斯丁主教的书信的\OriginalTerm{型式}{\GreekText{τύπον}}，他曾劝勉您的圣座（并非教导无知者，而是提醒知道者）要使那些他早已界定、如今认为宜于提醒你们的事，依照共同信仰的准则及天主教会，得以彻底执行。\par

菲尔穆斯，卡帕多细亚的凯撒勒雅主教说：至圣主教色勒斯丁的宗座和圣座先前已就此事给出了决定和\OriginalTerm{型式}{\GreekText{τύπον}}，通过致最蒙天主喜爱的诸位主教——即亚历山大的济利禄、耶路撒冷的犹文纳尔、得撒洛尼的鲁福斯，以及君士坦丁堡和安提约基雅的圣教会——的书信。我们也已遵循了这型式（既然为聂斯多略改正所定的期限早已过去，且自我们依最虔诚皇帝谕令抵达厄弗所城以来已过了许多时日，随后又耽搁了不少时间，以致皇帝所定的日期已过；且聂斯多略虽经传唤却未出庭），我们执行了这\OriginalTerm{型式}{\GreekText{τύπον}}，宣布了对他的合乎教规的判决。\par

阿卡狄乌斯，最可敬的主教及代表说：尽管我们的航行缓慢，且逆风尤其阻碍我们，以致我们不知能否如所期望的那样抵达目的地，然而因天主良好的眷顾……因此，我们渴望请求您的圣座，请您命令我们得知已由您的圣座所规定的事。\par

斐理伯，宗座代表及司铎说：我们向圣可敬的会议表示感谢，因为当你们聆听了我们圣蒙福教宗的书信之后，圣洁的成员们藉着我们（或你们）的圣洁声音，也藉着你们的圣洁欢呼与圣头联合。因为你们的圣座并非不知，整个信仰的头、宗徒们的头，是蒙福的宗徒伯多禄。如今，我们的卑微在饱受风浪、多次困扰之后终于抵达，我们请求你们命令将我们抵达之前在这圣会议中所行的事陈于我们面前，以便依照我们蒙福教宗和此次神圣集会的意见，我们也能批准他们的决定。\par

德奥多托，安塞拉主教说：全世界的天主已藉至宗教主色勒斯丁的书信和您的圣座的来临，显明了圣会议所宣判的正义。因为你们显明了至圣可敬主教色勒斯丁的热忱及其对虔诚信仰的关心。既然尊座极合理地渴望从关于罢免聂斯多略的会议记录中得知所行之事，尊座将完全确信判决的正义、圣会议的热忱，以及至虔诚圣主教色勒斯丁以洪亮声音所宣明的信仰的和谐；当然，在尊座完全确信之后，其余的事将加入当前的行动。\par

[\Emph{在《行传》中接着有两封塞莱斯定（Cœlestine）的短函，一封致皇帝，一封致济利禄（Cyril），但并未提及它们的内容，也未说明它们如何出现在这里，于是本次会议的记述便这样突然结束了。}]\par

% structure: p0111 source=h2 target=section reason=relative-heading-rank; base=h2

\section{第三次会议}

% structure: p0112 source=h3 target=subsection reason=relative-heading-rank; base=h2

\subsection{摘自《行传》}

耶路撒冷主教犹文纳（Juvenal）对最尊敬的阿尔卡丢（Arcadius）与普若耶克图（Projectus）二位主教以及最尊敬的长老斐理伯（Philip）说：昨天，当此神圣大公会议开会时，在尊驾在场的情况下，你们要求在宣读了至圣至福的大罗马主教塞莱斯定的信函后，宣读《行传》中关于异端者聂斯多略（Nestorius）被革职的记录。于是会议下令照办。恳请尊驾告知我们，你们是否已阅读过这些记录并理解其效力。\par

宗座代表斐理伯长老说：通过阅读《行传》，我们发现了贵圣会议关于聂斯多略所行之事的记载。我们从记录中发现，所有决定均符合教规与教会纪律。现在，我们也向尊驾请求，虽然可能多余，但贵会议所宣读的内容，如今应再次向我们宣读，以便我们能遵循至圣教宗塞莱斯定的\OriginalTerm{模式}{\GreekText{τύπῳ}}（他同样将此职责托付予我们）以及尊驾的意愿，并能够\OriginalTerm{确认}{\GreekText{βεβαιώσαι}}这项判决。\par

[\Emph{阿尔卡丢附议斐理伯的动议后，默农（Memnon）指示宣读记录，由首席公证员执行。}]\par

宗座代表斐理伯长老说：毫无疑问，事实上万世皆知，至圣至福的伯多禄（Peter），宗徒之\OriginalTerm{首领}{\GreekText{ἔξαρχος}}与元首，信德的柱石，以及天主教会（Catholic Church）的\OriginalTerm{根基}{\GreekText{θεμέλιος}}，从我们的主耶稣基督、人类救主与救赎者那里领受了天国的钥匙，并且赐予他束缚与赦免罪过的权能；他直至今日并永远活在他的继承人身上并审判他们。至圣至福的教宗塞莱斯定，按适当顺序，是他的继承者并占据其位，他派遣我们在此圣会议中代行其职；此会议是由最仁慈、最虔敬的皇帝们下令召集的，他们铭记并不断维护公教信仰。因为他们从他们最虔诚、最仁慈的祖先及圣德可念的先辈那里继承下来的宗徒训诲，直到如今都持守不渝，等等。\par

[\Emph{演讲中未再提及教宗特权。}]\par

最尊敬的阿尔卡丢主教、宗座代表说：聂斯多略使我们极为悲痛……既然他自愿使自己成为我们中的异己和流放者，我们遵循从起初由圣宗徒和天主教会（Catholic Church）传下来的规定（因为他们教导的是从我们的主耶稣基督所领受的），也遵循宗座至圣教宗塞莱斯定的\OriginalTerm{训令}{\GreekText{τύποις}}（他屈尊派遣我们作为他在此事务的执行者），并遵循圣会议的决议（我们给出此结论）：让聂斯多略知道，他被剥夺了一切主教尊位，与整个教会及所有司铎的共融隔绝。\par

罗马教会主教普若耶克图说：从阅读中极为清楚，等等……此外，我以圣宗座代表的职权，与我的弟兄们共同作为上述判决的\OriginalTerm{执行者}{\GreekText{ἐκβιβαστής}}，判定：前述聂斯多略是真理的敌人，信仰的败坏者，对其被指控之事有罪，应从主教尊位等级中被撤除，并进一步与所有正统司铎的共融隔绝。\par

亚历山大主教济利禄说：阿尔卡丢与普若耶克图二位至圣至敬的主教，以及罗马教会最虔诚的长老斐理伯所作的声明，在圣会议面前是显而易见的。因为他们是以宗座及整个西方虔敬神爱、至圣主教们的圣会议的名义作此声明。因此，应执行由至圣的、神爱的塞莱斯定主教所决定的事项，并普遍赞同在厄弗所（Ephesus）都会召开的圣会议对异端者聂斯多略所投的票；为此目的，应将昨天和今天的程序添加到已准备好的行传中，并呈交他们的圣座，以便他们按照惯例签署，表明他们与我们在教规上的一致。\par

最尊敬的阿尔卡丢主教、罗马教会代表说：根据此圣会议的记录，我们必须以我们的署名确认他们的教义。\par

圣会议说：既然最尊敬、最虔诚的阿尔卡丢与普若耶克图二位主教和代表，以及宗座代表斐理伯长老，已表明与我们同心，那么只需他们履行承诺，以签名确认会议记录，之后将会议记录出示给他们。\par

[\Emph{三人随后签署。}]\par

% structure: p0124 source=h2 target=section reason=relative-heading-rank; base=h2

\section{第七次会议（教规）}

遵照我们最虔诚的皇帝们之谕令，在厄弗所集合的神圣大公会议，致各省各城的主教、长老、执事及全体人民：\par

当我们按照皇帝陛下虔诚的谕令在厄弗所省会合时，有三十余人从我们中间退出，其分裂的领袖是安提约基雅主教若望。他们的名字如下：首先，上述叙利亚安提约基雅的若望，大马士革的若望，阿帕美亚的亚历山大，希拉波利斯的亚历山大，尼科美底亚的希美里乌斯，赫拉克利亚的弗里提拉斯，塔尔索的赫拉迪乌斯，阿纳扎尔布斯的马克西明，马尔基雅诺波利斯的狄奥多勒，特拉雅诺波利斯的伯多禄，厄默萨的保禄，赫拉克勒奥波利斯的波利克罗尼乌斯，提亚纳的欧提里乌斯，新凯撒利亚的默肋提乌斯，居鲁斯的狄奥多勒，加采东的阿普林吉乌斯，大劳迪西亚的玛加略，厄斯布斯的佐西斯，基里基雅科律库斯的撒路斯提乌斯，基里基雅卡斯塔巴拉的赫西基乌斯，穆特洛布拉卡的瓦伦提尼安，帕纳索斯的欧斯塔提乌斯，特奥多西亚的斐理伯，以及达尼尔、德克西安、犹利安、济利禄、奥林庇乌、狄奥革内、波利乌、费拉德尔菲亚的德奥凡、奥古斯塔的特拉扬、伊雷诺波利的奥勒良、阿拉杜斯的弥赛、托勒迈伊斯的赫拉迪乌斯。这些人并无基于司铎权威的教会共融特权，可借此伤害或有益于任何人；因为他们中有些人已被罢免，且因他们拒绝加入我们反对聂斯多略的法令，显然所有人都拥护聂斯多略和塞莱斯提乌斯的见解。神圣会议通过一项共同法令，将他们逐出一切教会共融，并剥夺他们一切司铎权能，使他们无法伤害或有益于任何人。\par

% structure: p0127 source=h3 target=subsection reason=relative-heading-rank; base=h2

\subsection{第一条}

鉴于那些因任何教会或个人原因被阻留而未能出席神圣会议、仍留在自己教区或城市的人，不应不知晓会议所颁布的事宜；因此我们通知您的圣德与爱德：若任何省区的都主教，弃绝神圣的大公会议而加入或日后加入叛教者的集会，或采纳或日后采纳塞莱斯提乌斯的道理，他绝不可在任何方面反对该省主教们，因为他已被逐出一切教会共融且不能行使其职务；他反而应在一切事上服从该省那些主教及邻近的正统都主教，并被剥夺其主教品级。\par

% structure: p0129 source=h3 target=subsection reason=relative-heading-rank; base=h2

\subsection{第二条}

若有省区主教未出席神圣会议，却加入或企图加入叛教者的集会；或在签署了聂斯多略的罢免后，又退回到叛教者的集会中；这些人根据神圣会议的法令，应被免除司铎职，剥夺其品级。\par

% structure: p0131 source=h3 target=subsection reason=relative-heading-rank; base=h2

\subsection{第三条}

若有城市或乡村的圣职人员，因正统信仰而被聂斯多略或其追随者禁止行使司铎职务，我们宣告他们应恢复原品级。此外，我们普遍禁止所有忠于正统大公会议的圣职人员，以任何方式服从那些已经或将要背教的教长。\par

% structure: p0133 source=h3 target=subsection reason=relative-heading-rank; base=h2

\subsection{第四条}

若有圣职人员背离，并公开或私下擅自主张聂斯多略或塞莱斯提乌斯的道理，神圣会议宣告这些人也应被罢免。\par

% structure: p0135 source=h3 target=subsection reason=relative-heading-rank; base=h2

\subsection{第五条}

若有任何人因恶行受到神圣会议或其本主教谴责，且聂斯多略（或其追随者）以其一贯的不加分辨，曾试图或日后试图不合教规地将这些人恢复至共融及原品级，我们宣告他们不得从中获益，但仍应被罢免。\par

% structure: p0137 source=h3 target=subsection reason=relative-heading-rank; base=h2

\subsection{第六条}

同样，若有任何人以任何方式试图推翻神圣会议在厄弗所就各案所颁的命令，神圣会议规定：若他们是主教或圣职人员，则完全丧失其职务；若为平信徒，则处以绝罚。\par

% structure: p0139 source=h3 target=subsection reason=relative-heading-rank; base=h2

\subsection{第七条}

当这些事项被宣读后，神圣会议规定：任何人不得提出、书写或编纂一个\OriginalTerm{不同的}{\GreekText{ἑτέραν}}的信经，以取代由圣神与圣父们在尼西亚所制定的信经。\par

但那些胆敢编纂不同信经，或将其引进或提供给那些愿意归向真理的认识者（无论他们来自异教、犹太教或任何异端）的人，若是主教或圣职人员，则主教被罢免主教职，圣职人员被开除圣职；若是平信徒，则予以绝罚。\par

同样，若发现有任何人（无论是主教、圣职人员或平信徒）持有或传授司铎卡里修斯所提出的《陈述》中关于天主独生子降生成人的道理，或以下所附的聂斯多略可憎不敬的道理，应受此神圣大公会议所定之罚。即：若为主教，应被撤职并贬黜；若为圣职人员，亦应开除圣职；若为平信徒，则被绝罚，如前所述。\par

% structure: p0143 source=h3 target=subsection reason=relative-heading-rank; base=h2

\subsection{第八条}

我们的弟兄、天主所爱的\Person{Rheginus}主教，以及他同样天主所爱的同僚、\Place{塞浦路斯}省的主教\Person{Zeno}和\Person{Evagrius}，向我们报告了一桩违反教会章程和\WorkTitle{圣宗徒法令}、并触及所有人自由的革新。因此，既然影响所有人的伤害需要更多关注，因为它们造成更大损害，尤其是当它们违反古老习惯时；况且这些曾向教务会议请愿的卓越人士已通过书面和口头告诉我们，\Place{安提约基亚}的主教就这样在\Place{塞浦路斯}举行了祝圣；因此，\Place{塞浦路斯}各圣教会的主教，应按照真福教父们的法令和古老习惯，不受争议或侵害地享有自行祝圣其卓越主教的权利。同样的规则应在各地其他教区和省份遵守，以便任何天主所爱的主教不得接管任何从最初开始就未曾属于他本人或其前任管辖的省份。但若有人强行夺取并征服[一省]，他应将其交出；以免违反教父们的法令；或以免在神圣职务的幌子下引入世俗荣誉的虚荣；或以免我们在不知不觉中逐渐丧失我们的主耶稣基督、众人的救赎主以其宝血赐予我们的自由。\par

因此，这神圣的大公会议决定，在每个省份中，从最初就属于它的权利，应根据古老通行习惯，保持不变且不受侵害地予以保全；每位都主教有权为自己的安全获取一份这些法令的副本。若有人提出与本决定相悖的规则，这神圣的大公会议一致决定其无效。\par

% structure: p0146 source=h3 target=subsection reason=relative-heading-rank; base=h2

\subsection{厄弗所同一神圣教务会议致\Place{潘菲利}省神圣教务会议关于曾任其都主教的\Person{Eustathius}的书信}

因为天主启示的圣经说：\Quote{凡事凭智慧而行}，那些接受司铎职务的人尤其应当仔细审查所有待处理的事务。因为对于那些如此度过一生的人，他们得以在属于他们的事上确立诚实的希望，并如顺风般被带往他们所渴望的事；因此，圣经的这句话确实很有道理。但有时，痛苦难忍的悲伤侵袭心灵，并残酷地蒙蔽它，以至于将人从追求必要之事中带走，并说服他相信本来有害的事是有益的。我们曾看到那位最卓越、最虔诚的主教\Person{Eustathius}遭受了类似的情况。因为证据表明他已被合法祝圣；但正如他所声称的，因某些人的极大搅扰，且进入了他未曾预料的环境，因此，尽管他完全有能力驳斥迫害者的诽谤，却因对事务异常不熟悉而拒绝与困扰他的困难争斗，并以某种我们不知道的方式递交了辞职书。然而，当他一旦受托司铎的牧职，就应当以属灵的热忱坚持它，并仿佛卸下一切去与困难斗争，且愉快地忍受他所承诺的辛劳。但既然他证明自己缺乏实际能力，与其说由于怯懦或懒惰，不如说由于不熟悉而遭遇此不幸，所以您的圣洁不得不任命我们最卓越、最虔诚的弟兄和同僚主教\Person{Theodore}为该教会的监督；因为让教会保持寡妇状态、让救主的羊群没有牧人地度日是不合理的。但当他来到我们这里哭泣时，并非与上述最虔诚的主教\Person{Theodore}争夺他的主教座或教会，而只是在此期间寻求主教的衔头和品秩，我们所有人都与这位老人一同悲伤，并视他的哭泣为己有，急忙查明上述[\Person{Eustathius}]是否受到合法罢免，或者他是否确实因某些人对他名誉散布无稽闲谈而提出的荒谬指控而被定罪。事实上我们得知并未发生此类事情，反而他的辞职被视为对上述\Person{Eustathius}的正式控告。因此，我们绝不责怪您的圣洁被迫祝圣上述最卓越的\Person{Theodore}主教接替其位。但是，既然过分反对此人的不切实际并不合适，而更应当怜悯这位年事已高、远离其出生之城和祖先居所的老年人，我们经司法裁决并毫无反对地宣布，他应保有其主教衔头和品秩以及共融。但条件如下：他不得祝圣，也不得凭个人权威接管并牧养一个教会；只有在被兄弟和同僚主教根据基督内的规定和爱德邀请作为助手或若偶然被任命时，方可如此。然而，如果您现在或将来为他决定更有利的安排，这也将使圣教务会议满意。\par

% structure: p0148 source=h3 target=subsection reason=relative-heading-rank; base=h2

\subsection{教务会议致教宗\Person{Celestine}的书信}

因天主的恩宠在首都\Place{厄弗所}集会的神圣教务会议，致最圣洁和我们的同工\Person{Cœlestine}，愿主赐予健康。\par

圣座对虔敬的热忱，以及对正确信德的关怀，令我们众人之救主天主喜悦感恩，堪受一切赞美。因为您在这等重大事务上的惯常做法是检验万事，并将教会的巩固视为己任。但既然一切所发生之事理当让圣座知晓，我们不得不写信（告知您）如下：藉着我们众人之救主基督的旨意，并依照最虔诚且爱基督的皇帝们的命令，我们从许多遥远分散的地区聚集在\Place{厄弗所}大城，总计超过二百位主教。然后，按照召集我们的爱基督的皇帝们的谕令，我们确定了神圣公会的会议日期为圣神降临节，全体一致同意，尤其因为皇帝的信函中规定：若有人未在指定时间到达，他便在良心上有所欠缺，在天主和人面前均无可推诿。最可敬的\Place{安提约基雅}主教\Person{若望}滞留未到；他并非出于诚心，也非因路途遥远造成阻碍，而是心中隐藏着那不讨天主喜悦的计谋和念头——这计谋和念头在他不久后到达\Place{厄弗所}时便显露无遗。\newline{}\par

因此，在指定的圣神降临节之后，我们将会议推迟了整整十六天；在此期间，许多主教和圣职人员染病，并为开销所累，甚至有人去世。这给伟大的公会带来了巨大损害，圣座不难察觉。因为他乖僻地利用如此漫长的拖延，使得许多从更远地区来的人反而先于他到达。\newline{}\par

然而十六天后，与他同行的一些主教中有两位都主教——\Place{阿帕梅雅}的\Person{亚历山大}和\Place{希拉波利}的\Person{亚历山大}——先于他到达。当我们多次抱怨最可敬的\Place{安提约基雅}主教\Person{若望}姗姗来迟时，我们被告知：\Quote{他命令我们向尊座宣布，若有什么事情耽搁他，不要推迟公会，而要做正当的事。}接到这一消息后——既然从他的拖延以及刚才向我们宣布的情况来看，显然他拒绝出席公会，或是出于对\Person{聂斯多略}的友谊，或是因为他曾受其管辖的教会中的圣职人员，或是顾及某些人为他所作的请求——神圣公会便在\Place{厄弗所}那座以\Person{玛利亚}命名的大教堂中召开了会议。\newline{}\par

但当众人热忱地聚集时，只有\Person{聂斯多略}一人缺席公会，于是神圣公会依照教规，通过主教们向他发出了第一次、第二次和第三次劝告。但他以士兵包围自己的住所，对抗教会法律，既不露面，也不为他邪恶的亵渎言论作出任何赔偿。\newline{}\par

此后，宣读了\Place{亚历山大}教会的至圣至敬的主教\Person{济利禄}写给他的信件，神圣公会批准这些信件为正统且无可指摘（\OriginalTerm{正确且无可指摘}{\GreekText{ὀρθῶς} \GreekText{καὶ} \GreekText{ἀλήπτως} \GreekText{ἔχειν}}），并且与天主所默感的圣经以及由圣教父们（从前在\Place{比提尼亚}的\Place{尼西亚}召开的大公会）所传承并阐述的信德毫无相悖之处，正如圣座您在仔细审查后也予以证实的那样。\newline{}\par

另一方面，宣读了\Person{聂斯多略}写给上述我们最可敬的同僚兼同工\Person{济利禄}的信件，神圣公会认为信件中所教导的内容完全背离宗徒及福音的信德，充满许多离奇的亵渎之词。\newline{}\par

同样宣读了其最不虔敬的释经著作，以及圣座您写给他的信函，他在信中被正当地谴责为写下亵渎之言，并在其私人释经中插入不虔敬的观点（\OriginalTerm{言论}{\GreekText{φωνᾶς}}）。此后，对他宣告了公正的免职判决；这一判决尤其公正，因为他如此远离悔改，或是对其亵渎之事作任何承认，尽管他拥有\Place{君士坦丁堡}教会，甚至在\Place{厄弗所}大城中，他还向某些都主教（他们并非无知，而是博学且敬畏天主）发表了一篇讲道，在其中他竟敢说：\Quote{我不承认一个两三个月大的天主，} 还说了比这更过分的话。\newline{}\par

因此，作为一种亵渎且极其有害的异端，它颠倒我们最纯洁的宗教（\OriginalTerm{宗教}{\GreekText{θρησκείαν}}），并从根本上推翻整个奥迹（即降生成人）的安排，我们将其推翻，如上所述。但看来，那些拥有基督真诚之爱并在主内热忱的人不可能不经历许多考验。因为我们曾希望最可敬的\Place{安提约基雅}主教\Person{若望}会赞扬公会的勤勉和虔诚，或许会责备\Person{聂斯多略}被免职的迟缓。但一切与我们的期望相反。因为他被发现是神圣公会的敌人，甚至是教会正统信仰的最具敌意的敌人，正如这些事所表明的。\newline{}\par

因为他一来到厄弗所，甚至尚未抖落旅途的尘土，或更换旅行服装，就召集了那些支持奈斯多利并发出亵渎反对他们首领的人，几乎是对基督的荣耀加以嘲弄，并且聚集了一伙，我想，大约三十人，具有主教的名号（其中一些没有教区，四处游荡，没有辖区；另一些则因严重过失多年以前就被剥夺了都城主教职位；还有白拉奇派和塞莱斯提乌的追随者，以及一些被逐出特萨利亚的人），就竟敢犯下无人曾行过的恶行。因为他独自起草了一份他称之为\Quote{革职令}的文件，辱骂和责备至圣可敬的亚历山大主教济利禄和至可敬的厄弗所主教默农、我们的弟兄和同工，我们中无人知情，甚至那些被辱骂的人也不知道正在发生什么，也不清楚他们凭什么胆敢如此行。但是他们漠视天主对这种行为的义怒，不顾教会法规，忘记了他们正在如此行事奔向毁灭，以绝罚的名义，他们辱骂了整个主教会议。然后将他们的这些行为张贴在公共布告栏上，任凭愿意阅读的人阅读，张贴在剧院外面，好让他们把自己的不敬奉为表演。但这还不是他们胆量的极限；反而好像他们做了合乎教规的事，居然敢把他们所做的传到至虔敬并爱基督的皇帝们耳中。情况如此，至圣可敬的亚历山大主教济利禄和至可敬的厄弗所城主教默农，提交了他们自己编写的一些文书，控告至可敬的若望主教和那些与他一同行了这事的人，恳求我们神圣的主教会议按照教规传唤若望和他的人，以便他们为它们的胆大妄为道歉，并且如果他们有任何控诉，可以发言并证明；因为在他们的书面革职令，或者不如说是辱骂单中，他们以此为借口声称：\Quote{他们是阿波利纳里派、亚略派和欧诺米派，因此被我们罢免。}因此，当那些遭受他们辱骂的人在场时，我们再次必然聚集在大教堂中，共有二百多位主教，并在两天内通过第一次、第二次和第三次传唤，召叫若望及其同伴到主教会议来，以便他们审查那些被辱骂的人，并作解释，告诉他们是什么原因促使他们起草了革职判决；但他不敢前来。\par

但是，如果他真能证明前述圣人们是异端者，他理当前来，并证明那项被视为真实且无可置疑的罪行，从而导致对他们作出轻率的判决。但他被自己的良心定罪，没有前来。他所图谋的是这样：因为他认为当那毫无根据且最不义的辱骂被消除后，主教会议针对异端者奈斯多利所作的正确表决也会随之失效。因此，我们理当愤怒，决定依法对他和他的人施加同样的惩罚，即他违法地加诸那些未被证明有罪的人身上的惩罚。但是，尽管最正义且依法他本应遭受此罚，但希望因我们的忍耐而克服他的鲁莽，我们将此保留给圣座裁决。与此同时，我们剥夺了他们的共融，并取走了他们所有司铎的权能，使他们无法凭借其见解作恶。因为那些如此凶猛、残忍、不合教规地惯于冲向如此可怕和最邪恶之事的人，岂不应该被剥夺他们实际上并不拥有的作恶的权能吗？\par

至于我们的弟兄和同工，即主教济利禄和默农，他们曾忍受他们的责备，我们都与他们共融，在[控告者]的鲁莽之后，我们共同举行并正在进行礼仪，一起举行圣祭，使他们的书面把戏无效，从而表明它完全没有效力和作用。因为这不过是辱骂而非其他。因为三十个人，其中一些带有异端的印记，一些没有教区并被[从他们的教区]驱逐，能组成什么样的主教会议呢？或者，它能有什么力量来对抗从全世界聚集的主教会议呢？因为与我们同坐的有至可敬的阿尔卡狄乌斯主教和普罗耶克图斯主教，以及至圣的司铎斐理伯，他们都是圣座派遣的，他们向我们传达了圣座之临在，并充满\OriginalTerm{宗座}{\GreekText{τῆς} \GreekText{ἀποστολικῆς} \GreekText{καθέδρας}}的地位。那么，愿圣座为此事发怒。但是，如果允许那些愿意向较大教区倾泻辱骂的人，并且不合法、不合教规地下判决，或者更确切地说，向那些他们无权管辖的人发出辱骂，向那些因信仰而忍受如此巨大争战的人发出辱骂——正因如此，如今因圣座的祈祷，虔诚得以闪耀——[我要说，如果这一切被容忍]，教会事务将陷入最大的混乱。但当那些胆敢如此行的人得到适当惩戒时，一切骚动将停止，对教规应有的尊敬将为所有人所遵守。\par

当神圣公会议听取了对最不虔敬的佩拉奇派和塞莱斯廷派、即塞莱斯提乌斯、佩拉奇乌斯、尤利安、普雷西迪乌斯、弗洛鲁斯、马尔塞利安、奥朗提乌斯以及倾向类似错误者所作的革除处分后，我们也认为正确（\OriginalTerm{认为正确}{\GreekText{ἐδικαιώσαμεν}}）应使贵圣座关于他们的决定保持稳固且坚定。我们全体同心，视他们为已被革除。为使贵方完全明了所行一切，我们已将法令文本及公会议签署名单副本寄送。祈求主，深爱且极渴望的诸位，在主内坚强并记念我们。\par

% structure: p0162 source=h3 target=subsection reason=relative-heading-rank; base=h2

\subsection{反对不虔敬的墨萨利派（亦称尤赫特派和恩图西亚斯特派）之定义}

当最虔诚且最热心的主教瓦莱里安和安菲洛奇乌斯来到我们这里时，他们提议我们共同考虑墨萨利派人，即尤赫特派或恩图西亚斯特派的情况，这些人在潘菲利亚盛行，或者无论这最污秽的异端被称作其他什么名字。当我们思考此事时，最虔诚且最热心的主教瓦莱里安向我们提交了一份在蒙福的西辛尼乌斯治下于伟大的君士坦丁堡城针对他们拟定的公会议章程。我们所读到的内容得到众人赞赏，认为撰写恰当且案情陈述得当。我们全体——包括最虔诚的主教瓦莱里安和安菲洛奇乌斯以及潘菲利亚和吕考尼亚全省所有最虔诚的主教——都认为该公会议章程中所包含的一切应被确认且不可撤销；同时，在亚历山大所采取的行动也应被确定，以便全省范围内所有属于墨萨利派或恩图西亚斯特派异端、或被怀疑感染该异端的人，无论是圣职人员还是平信徒，都应聚集；如果他们按照上述公会议所颁布的法令以书面方式诅咒自己的错误，那么，若是圣职人员，可保留其职务；若是平信徒，则允许其共融。但如果他们拒绝诅咒，那么，若是司铎、执事或任何其他教会品级，应将其从圣职和其品级中驱逐，并排除于共融之外；若是平信徒，则应予以绝罚。\par

此外，凡被证实犯有该异端者，不再被允许管理我们的隐修院，以免莠子被撒播并增长。我们命令最虔诚的主教瓦莱里安和安菲洛奇乌斯以及全省其他最可敬的主教们注意执行此法令。除此之外，还认为应诅咒该异端的污秽之书，名为\WorkTitle{苦修集}，因其由异端者编纂；最虔诚且热心的瓦莱里安带来了该书的副本。同样，凡在其人民中间发现的任何带有他们不虔敬气味的东西，都应受诅咒。\par

此外，当他们聚集时，应由他们以书面方式提交那些对和谐、共融及安排（\OriginalTerm{安排或管理}{dispositionem vel dispensationem}）有用且必要的事项。但若在当前事务中出现任何问题，且该问题被证明是困难和含糊的，凡未经最虔诚的主教瓦莱里安和安菲洛奇乌斯以及全省其他主教认可的事项，他们应参照所写的来讨论一切。若吕西亚或吕考尼亚的最虔诚主教被忽略，无论如何，无论他属于哪个省，都不得遗漏一位都主教。将这些内容记入法令，以便若有需要者，能找到如何更详尽地向他人解释这些事。\par

% structure: p0166 source=h3 target=subsection reason=relative-heading-rank; base=h2

\subsection{公会议关于尤普雷皮乌斯和济利禄事宜的法令}

最虔诚的主教尤普雷皮乌斯和济利禄的请求——在他们提交的文件中阐明的——是正当的。因此，根据神圣教规和从古代习惯中获得法律效力的世俗法律，不得在欧罗巴诸城进行任何革新，而应按古代习惯由历来管辖它们的主教管理。因为既然从未有过都主教拥有其他权力，所以今后也不得背离古代习惯。\par

两位色雷斯主教，即比扎（比齐亚）的尤普雷皮乌斯和科埃莱的济利禄，引发了一项法令，请求保护他们免受其都主教——已倒向安提约基雅的若望一派的赫拉克勒亚的弗里提拉斯——的侵害，同时请求确认先前同时担任两个主教座堂的做法。公会议批准了这两项请求。\par

\SourceRecord{Translated by Henry Percival.；Nicene and Post-Nicene Fathers, Second Series；Vol. 14.；Edited by Philip Schaff and Henry Wace.；Buffalo, NY: Christian Literature Publishing Co.,；1900.；<http://www.newadvent.org/fathers/3810.htm>.}
